\begin{frame}[fragile]
	\titlepage
\end{frame}

\begin{latexonly}
section{Activit�s} % (fold)
\label{sec:activit�s}

	\begin{frame}[fragile]
		\frametitle{Sur cahier d'exercices}
		\begin{center}
				\includegraphics[width=.9\textwidth]{RC_activite01.png}
		\end{center}
	\end{frame}

	\begin{frame}[fragile]
		\frametitle{Sur cahier d'exercices}
		\begin{center}
				\includegraphics[width=.9\textwidth]{RC_activite02.png}
		\end{center}
	\end{frame}

	\begin{frame}[fragile]
		\frametitle{Sur cahier d'exercices}
		\begin{center}
				\includegraphics[width=.9\textwidth]{RC_activite03.png}
		\end{center}
	\end{frame}

	\begin{frame}[fragile]
		\frametitle{Sur cahier d'exercices}
		\begin{center}
				\includegraphics[width=.9\textwidth]{RC_activite04.png}
		\end{center}
	\end{frame}
% section activit�s (end)
\end{latexonly}



\section{D�finition} % (fold)
\label{sec:d�finition}
\begin{frame}[fragile]
	\begin{defi}[Racine carr�e]
	    Soit $a$ un nombre positif, la racine carr�e de $a$ est le nombre
	    positif qui �lev� au carr� donne $a$.\par \pause
	    On le note $\sqrt{a}$ et on a~: $(\sqrt{a})^2=a$.
	\end{defi}
\end{frame}

\begin{frame}[fragile]

Soit $a$ un nombre. Que vaut $\sqrt{a^2}$~? \pause
\begin{enumerate}
	\item  $\sqrt{a^2}=a$, si $a$ est positif ;
	\item  $\sqrt{a^2}=-a$, si $a$ est n�gatif.
\end{enumerate}

\begin{eg}
	\begin{latexonly}
		\begin{itemize}
			\item Soit $a=5$. \pause Le nombre est positif, on a $\sqrt{5^2}=5$.\pause
			\item Soit $a=-6$. \pause Le nombre est n�gatif, on a $\sqrt{-6)^2}=\pause -(-6)=6$.
		\end{itemize}
	\end{latexonly}

	\begin{wimsonly}
		\def{integer k=random(-1,1)}
		\def{integer a= random(2..20)}
		\def{integer ka=\a*\k}

		\if{\k=1}{\def{text signe=positif}}{\def{text signe=n�gatif}}

		\if{\k=1}{\def{text solution= \ka}}{\def{text solution= $-(\ka)=\a$}}

		On cherche � calculer $A=\sqrt{(\ka)^2}$.<br> On remarque que $\ka$ est \signe, donc on a $A=$ \solution.
	\end{wimsonly}
\end{eg}
	\begin{rmq}
		On remarque que grace au carr�, on prend toujours la racine carr�e d'un nombre positif, et que le r�sultat final est �galement un nombre positif, comme cela �tait dit dans la d�finition.
	\end{rmq}
\pause
\end{frame}

\begin{frame}[fragile]
\begin{rmq}
	Ne pas confondre avec $(\sqrt{a})^2$, qui n'existe que si $a$ est positif et alors on a~:$(\sqrt{a})^2=a$.
\end{rmq}

\pause
	\begin{latexonly}
		\begin{eg}
			On a bien $\sqrt{5}^2=5$ et $\sqrt{-7}^2$ \pause qui n'existe pas !
		\end{eg}
	\end{latexonly}

	\begin{wimsonly}
		\begin{eg}[\reload{<img src = "gifs/doc/etoile.gif" alt = "rechargez" width = "20" height = "20">}]
			\def{integer a=random(2..20)}
			\def{integer b=random(2..20)}
			On a bien
			$\sqrt{\a}^2=\a$ et $\sqrt{-\b}^2$ qui n'a pas de sens.
		\end{eg}
	\end{wimsonly}

\begin{enumerate}
	\item  Dans l'expression $\sqrt{a}$, $a$ est le radicande~;\pause
    \item  On a toujours $\sqrt{a} \geq 0$, voir la d�finition~;\pause
    \item  Les nombres dont la racine carr� est un entier sont appel�s \emph{carr�s parfaits}.\par
\end{enumerate}

\pause
\begin{latexonly}
	\begin{eg} $\sqrt{25}=5$, on dit alors que 25 est un carr� parfait, car sa racine carr�e est un entier.
	\end{eg}
\end{latexonly}

\begin{wimsonly}
	\def{integer a=random(2..20)}
	\def{integer aa= \a*\a}
\begin{eg}[Calculons l'expression \( A\) \reload{<img src = "gifs/doc/etoile.gif" alt = "rechargez" width = "20" height = "20">}	]
		<div class="wimscenter" >
		\(	\begin{split}
   			A & = \sqrt{\aa} \\
   			A & =\a
   		\end{split} \)
		</div>
	On dit que \( \aa \) est un carr� parfait, car sa racine carr�e est un entier.
\end{eg}
\end{wimsonly}

\end{frame}
% section d�finition (end)



\section{Propri�t�s} % (fold)
\label{sec:propri�t�s}

\begin{frame}[fragile]
	\begin{prop}
	    Soit $a$ et $b$ deux nombres positifs, on a~:
	    \[ \sqrt{a}\times \sqrt{b}=\sqrt{ab}\] \pause
	    De plus, si $b \not = 0$, on a aussi~:
	    \[ \sqrt{\frac{a}{b}}=\frac{\sqrt{a}}{\sqrt{b}} \]
	\end{prop}
\pause


\begin{latexonly}
	\begin{eg}
		\begin{columns}
			\begin{column}{0.45\textwidth}
			\begin{eqnarray*}
				A & = & \sqrt{18} \times \sqrt{2} \\ \pause
			    A & = & \sqrt{36}   \\ \pause
			    A & = & 6 \pause
			\end{eqnarray*}
			\end{column}


			\begin{column}{0.45\textwidth}
		\begin{eqnarray*}
			B & = & \sqrt{\frac{9}{4}} \\ \pause
		    B & = &\frac{\sqrt{9}}{\sqrt{4}}   \\ \pause
		    B & = & \frac{3}{2}
		\end{eqnarray*}
			\end{column}
		\end{columns}
	\end{eg}
\end{latexonly}

\begin{wimsonly}
	\begin{eg}[\reload{<img src = "gifs/doc/etoile.gif" alt = "rechargez" width = "20" height = "20">}]
		\begin{columns}
			\begin{column}{0.45\textwidth}
				\def{matrix a=2,8,16,4
				2,18,36,6
				2,32,64,8
				2,50,100,10
				5,20,100,10
				10,10,100,10
				2,72,144,12
				4,36,144,12
				9,16,144,12
				}
				\def{text a=randomrow(\a)}
				\def{integer a1=\a[1]}
				\def{integer a2=\a[2]}
				\def{integer a3=\a[3]}
			\begin{eqnarray*}
				A & = & \sqrt{\a1} \times \sqrt{\a2} \\ \pause
			    A & = & \sqrt{\a3}   \\ \pause
			    A & = & \a[4] \pause
			\end{eqnarray*}
			\end{column}


			\begin{column}{0.45\textwidth}
				\def{integer bb1=random(2..15)}
				\def{integer b1=(\bb1)^2}
				\def{integer bb2=random(2..15)}
				\def{integer b2=(\bb2)^2}
				\def{rational r=\bb1 / \bb2}
				\def{text comp1=\bb1/\bb2}
				\def{text comp2=\r}
				\if{\comp1=\comp2}{\def{text sol=}}{\def{text sol=B&=&\r}}
		\begin{eqnarray*}
			B & = & \sqrt{\frac{\b1}{\b2}} \\ \pause
		    B & = &\frac{\sqrt{\b1}}{\sqrt{\b2}}   \\ \pause
		    B & = & \frac{\bb1}{\bb2} \\
			\sol
		\end{eqnarray*}
			\end{column}
		\end{columns}
	\end{eg}
\end{wimsonly}
\end{frame}



\begin{frame}[fragile]
	\frametitle{Attention}
	\begin{latexonly}
	   	\begin{columns}
	   		\begin{column}{0.45\textwidth}
	   		\begin{eqnarray*}
	   		    A & = & \sqrt{9}+\sqrt{16} \\ \pause
	   		    A & = & 3+4  \\ \pause
	   		    A & = & 7
	   		\end{eqnarray*}
	   		\end{column}
	    \pause

	   		\begin{column}{0.45\textwidth}
	   		\begin{eqnarray*}
	   			B & = & \sqrt{9+16} \\ \pause
	   			B & = & \sqrt{25} \\ \pause
	   			B & = & 5
	   		\end{eqnarray*}
	   		\end{column}
	   	\end{columns}
	\end{latexonly}

\begin{wimsonly}
   	\begin{columns}
   		\begin{column}{0.45\textwidth}
			\def{matrix list=3,4,5
			5,12,13
			7,24,25
			9,40,41}
			\def{text a=randomrow(\list)}
			\def{integer A1=\a[1]^2}
			\def{integer A2=\a[2]^2}
			\def{integer a3=\a[1]+\a[2]}
			\def{integer A3=\A1+\A2}
   		\begin{eqnarray*}
   		    A & = & \sqrt{\A1}+\sqrt{\A2} \\ \pause
   		    A & = & \a[1]+\a[2]  \\ \pause
   		    A & = & \a3
   		\end{eqnarray*}
   		\end{column}
    \pause

   		\begin{column}{0.45\textwidth}
   		\begin{eqnarray*}
   			B & = & \sqrt{\A1+\A2} \\ \pause
   			B & = & \sqrt{\A3} \\ \pause
   			B & = & \a[3]
   		\end{eqnarray*}
   		\end{column}
   	\end{columns}
\end{wimsonly}
	\pause
	\begin{rmq}
	    Soit $a$ et $b$ des nombres strictement positifs, on a toujours~:
	    \[ \sqrt{a}+\sqrt{b}\not=\sqrt{a+b} \]
	\end{rmq}
\end{frame}


\begin{frame}[fragile]
On �crira toujours une racine carr�e sous sa forme la plus simple
possible.
\pause

\begin{latexonly}
	\begin{exo}
	    �crire les nombres donn�s sous la forme la plus simple.
	    $A=\sqrt{20}$~; $B=\sqrt{125}$~; $C=\sqrt{27}-3\sqrt{75}$
	\end{exo}


	\begin{sol}
	     On a $A=\sqrt{20} =$ \pause $\sqrt{2^2\times 5}=$\pause$ 2\sqrt{5}$~; \par
		 $B=\sqrt{125}= $\pause$ \sqrt{5^2\times 5}=$\pause$ 5\sqrt{5}$ ; \par
		 $C=\sqrt{27}-3\sqrt{75}= $\pause$ \sqrt{3^2\times 3}-3\sqrt{5^2\times 3}=$\pause$ 3\sqrt{3}-15\sqrt{3}=-12\sqrt{3}$
	\end{sol}
\end{latexonly}

\begin{wimsonly}



	\def{integer a=random(2..16)}
	\def{integer a2=\a^2}
	\def{matrix primes=2, 3, 5, 7, 11}
	\def{matrix list=shuffle(\primes)}
	\def{integer k=\list[1]}
	\def{integer n=\k*\a2}

	\begin{eg}[ \reload{<img src = "gifs/doc/etoile.gif" alt = "rechargez" width = "20" height = "20">}	]
	    �crire le nombre donn� sous la forme la plus simple : \( A=\sqrt{\n}\)
		<br>
		<br>
		On a :
		\begin{eqnarray*}
			A & = & \sqrt{\n} \\
			A & = & \sqrt{\k \times \a2} \\
			A & = & \sqrt{\k \times \a^2} \\
			A & = & \sqrt{\a^2} \times \sqrt{\k} \\
			A & = & \a\sqrt{\k}
		\end{eqnarray*}
	\end{eg}
\end{wimsonly}
\end{frame}

\begin{frame}[fragile]
	\begin{exo}
		\begin{itemize}
			\item \exercise{module=adm/raw&job=lightpopup&emod=H3/algebra/oefracine.fr&parm=cmd=new;exo=arrondie1;&option=noabout}{Calculer une racine carr�e (utilisation de la calculatrice)}
			\item \exercise{module=adm/raw&job=lightpopup&emod=H3/algebra/oefracine.fr&parm=cmd=new;exo=corres5;&option=noabout}{Simplifier des racines carr�es - Correspondance}
			\item \exercise{module=H3/algebra/oefracine.fr&cmd=new&exo=ecriture0&exo=ecriture&qnum=1&scoredelay=&qcmlevel=3;&option=noabout}{�criture r�duite}
			\item \exercise{module=H3/algebra/oefracine.fr&cmd=new&exo=expression&qnum=1&scoredelay=&qcmlevel=3;&option=noabout}{Calcul d'expression}
			\item \exercise{module=H3/algebra/oefracine.fr&cmd=new&exo=liaison5&qnum=1&scoredelay=&qcmlevel=3;&option=noabout}{Racines et nombres - Correspondance}
		\end{itemize}
	\end{exo}
\end{frame}
% section propri�t�s (end)

\end{document}

